% =====================================================================
%  CHAPTER 13: 食品安全监督管理
% =====================================================================
\chapter{食品安全监督管理}
\setcounter{chapter}{13}
\chapterintro{
\keyword{掌握}GMP的概念；掌握HACCP的概念和7个基本原则。\keyword{熟悉}食品安全法。
}

% =====================================================================
%  第一节 GMP（良好生产规范）
% =====================================================================
\section{GMP（良好生产规范）}

\begin{keybox}
\textbf{GMP（Good Manufacturing Practice）：}良好生产规范，是保证食品具有高度安全性的良好生产管理体系。GMP是注重生产过程中产品质量与卫生安全的自主性管理制度。

\textbf{GMP的四大管理要素（4M）：}
\begin{enumerate}[leftmargin=*]
    \item \textbf{Man（人员）：}配备合适且能胜任的从业人员，健康检查合格，经过食品安全知识培训。
    \item \textbf{Material（物料）：}选用良好的原料，原料验收标准，供应商管理。
    \item \textbf{Machine（设备）：}采用标准规范的厂房设施和机械设备，厂房布局合理（人流物流分离），设备易于清洗消毒。
    \item \textbf{Method（方法）：}遵照既定的最优方法进行生产，标准操作规程（SOP），工艺参数控制。
\end{enumerate}

\textbf{GMP的基本精神：}
\begin{enumerate}[leftmargin=*]
    \item 降低生产过程中的人为错误。
    \item 防止食品在生产过程中遭到污染和品质劣变。
    \item 建立企业自主性质量保证体系。
\end{enumerate}
\end{keybox}

% =====================================================================
%  第二节 HACCP（危害分析与关键控制点）
% =====================================================================
\section{HACCP（危害分析与关键控制点）}

\subsection{HACCP概念}

\begin{keybox}
\textbf{HACCP（Hazard Analysis and Critical Control Point）：}危害分析与关键控制点。是"从农田到餐桌"全过程的食品安全预防体系。通过对食品生产全过程的危害分析，确定关键控制点并采取控制措施，将危害降低到可接受水平。

\textbf{HACCP建立的前提条件：}HACCP体系必须建立在GMP和SSOP的基础之上。
\end{keybox}

\subsection{HACCP的7个基本原则}

\begin{keybox}
\textbf{HACCP的7个基本原则（掌握*）：}

\begin{enumerate}[leftmargin=*]
    \item \textbf{进行危害分析（Hazard Analysis）：}对食品生产全过程进行系统的危害分析，识别所有潜在的生物性、化学性和物理性危害。对每个危害进行风险评估。

    \item \textbf{确定关键控制点（Critical Control Point, CCP）：}确定能够实施控制的、对防止和消除食品安全危害或将其降低到可接受水平所必需的某一步骤。CCP的三种控制形式：防止危害发生、消除危害、降低到可接受水平。

    \item \textbf{建立关键限值（Critical Limit, CL）：}在CCP上设定的最大值和/或最小值，用于将危害控制在一定范围内。如巴氏消毒最低72$^\circ$C、15秒。

    \item \textbf{建立CCP的监控系统（Monitoring）：}按照计划对CCP进行观察和测量，以判断CCP是否处于控制之中。监控四大要素：监控什么（What）、怎样监控（How）、监控频次（Frequency）、谁监控（Who）。

    \item \textbf{建立纠偏措施（Corrective Action）：}当监控发现CCP偏离关键限值时，立即采取的措施。包括隔离和评估受影响产品、纠正偏离原因、防止再次发生、记录全部纠偏活动。

    \item \textbf{建立验证程序（Verification）：}确认HACCP体系是否有效运行。验证的三个维度：适宜性、一致性、有效性。

    \item \textbf{建立记录保持程序（Record-keeping）：}所有HACCP相关活动必须有完整、准确的记录。文件包括：HACCP计划文本、危害分析表、CCP监控记录、纠偏记录、验证记录。
\end{enumerate}
\end{keybox}

% =====================================================================
%  第三节 食品安全法
% =====================================================================
\section{食品安全法}

\begin{defbox}
\textbf{食品安全法（2015年修订）：}确立了以食品安全风险监测和评估为基础的科学管理制度。2015年10月1日修订版正式实施，被称为"史上最严"食品安全法。

\textbf{四大原则：}预防为主、风险管理、全程控制、社会共治。

\textbf{核心制度：}食品安全风险监测制度、食品安全风险评估制度、食品安全标准制度、食品生产经营许可制度、食品召回制度、食品安全信息发布制度。
\end{defbox}

% =====================================================================
%  复习题
% =====================================================================
\reviewcontent{
\section{复习题}

\subsection*{一、名词解释}

\begin{enumerate}[leftmargin=*, itemsep=6pt]
    \item \textbf{GMP（Good Manufacturing Practice）}
    \item \textbf{HACCP（Hazard Analysis and Critical Control Point）}
    \item \textbf{CCP（关键控制点）}
    \item \textbf{CL（关键限值）}
\end{enumerate}

\subsection*{二、简答题}

\begin{enumerate}[leftmargin=*, itemsep=8pt]
    \item \textbf{【重点】}简述GMP的概念和四大管理要素（4M）。
    \item \textbf{【重点】}简述HACCP的定义及7个基本原则。
    \item \textbf{【重点】}简述HACCP的7个基本原则中每个原则的主要内容。
    \item 简述食品安全法（2015年修订）的基本特征和四大原则。
\end{enumerate}

\subsection*{参考答案要点}

\subsubsection*{名词解释（要点）}

\begin{enumerate}[leftmargin=*, itemsep=4pt]
    \item \textbf{GMP：}良好生产规范，是保证食品具有高度安全性的良好生产管理体系。
    \item \textbf{HACCP：}危害分析与关键控制点，是"从农田到餐桌"全过程的食品安全预防体系。
    \item \textbf{CCP：}关键控制点，能够实施控制的、对防止和消除食品安全危害或将其降低到可接受水平所必需的某一步骤。
    \item \textbf{CL：}关键限值，在CCP上设定的最大值和/或最小值，用于将危害控制在一定范围内。
\end{enumerate}

\subsubsection*{简答题（要点）}

\begin{enumerate}[leftmargin=*, itemsep=6pt]
    \item \textbf{GMP：}良好生产规范，保证食品具有高度安全性的良好生产管理体系。\textbf{4M：}Man（人员——合适且胜任的从业人员）、Material（物料——选用良好的原料）、Machine（设备——标准规范的厂房设施和机械设备）、Method（方法——遵照既定的最优方法进行生产）。
    \item \textbf{HACCP定义：}危害分析与关键控制点，"从农田到餐桌"全过程的食品安全预防体系。\textbf{7原则：}①进行危害分析；②确定关键控制点（CCP）；③建立关键限值；④建立CCP的监控系统；⑤建立纠偏措施；⑥建立验证程序；⑦建立记录保持程序。
    \item \textbf{7原则内容：}①危害分析——识别生物、化学、物理性危害；②确定CCP——防止危害发生、消除危害、降低到可接受水平；③建立CL——最大值和/或最小值；④监控——What/How/Frequency/Who；⑤纠偏——立即隔离评估、纠正原因、防止再发、记录；⑥验证——适宜性、一致性、有效性；⑦记录保持——完整准确的文档记录。
    \item \textbf{食品安全法（2015年修订）：}确立以食品安全风险监测和评估为基础的科学管理制度。\textbf{四大原则：}预防为主、风险管理、全程控制、社会共治。
\end{enumerate}

\newpage
}
